%-------------------------------------------------------------------------------
% OTUS Run E / Tier A working paper
% Compile with: pdflatex main && bibtex main && pdflatex main && pdflatex main
%
% This is a deliberately structured skeleton.  Tables marked [RUN D] contain
% measured repository results; all [TODO] markers are for the Run E numbers.
%-------------------------------------------------------------------------------
\documentclass[11pt,twocolumn]{article}

\usepackage[margin=2.2cm]{geometry}
\usepackage{amsmath,amssymb}
\usepackage{booktabs}
\usepackage{xcolor}
\usepackage{graphicx}
\usepackage{hyperref}
\usepackage[numbers,sort&compress]{natbib}


\newcommand{\pt}{\ensuremath{p_{\mathrm{T}}}}
\newcommand{\ptvec}{\ensuremath{\vec{p}_{\mathrm{T}}}}
\newcommand{\zspace}{\ensuremath{\mathcal{Z}}}
\newcommand{\xspace}{\ensuremath{\mathcal{X}}}
\newcommand{\E}{\ensuremath{\mathcal{E}}}
\newcommand{\D}{\ensuremath{\mathcal{D}}}
\newcommand{\todo}[1]{\textcolor{red}{\textbf{[TODO: #1]}}}

\title{A Physics-Cost Neural-OT Upgrade of the Sliced-Wasserstein Autoencoder
for Detector-Level $J/\psi\to\mu^+\mu^-$ Simulation}

\author{[Author list -- TODO]}

\date{\today}

\begin{document}
\maketitle

\begin{abstract}
The OTUS sliced-Wasserstein autoencoder learns a stochastic detector response
between truth-level dimuon four-vectors and CMS detector-level four-vectors
without per-event paired simulation.  The literal 2021 objective matches broad
coordinate distributions but misses the $\sim30~\mathrm{MeV}$ $J/\psi$ mass shell.
This paper reports two incremental upgrades: (i) a Run-D baseline adding
pair-level sliced Wasserstein and a relative-mass cycle constraint, which
reduces the simulation mass Wasserstein-1 from $0.95~\mathrm{GeV}$ to
$0.026~\mathrm{GeV}$; and (ii) a Run-E ``Tier A'' upgrade that replaces blind random
slicing with a physics-cost Sinkhorn divergence and adversarially learned
slicing directions, parameterizes the transport in cylindrical coordinates as a
short stochastic residual flow, and selects checkpoints with hard mass-shape
gates.  Modern classifier two-sample, coverage, and per-event stochasticity
diagnostics are introduced in place of marginal histogram checks alone.
\todo{Abstract final numbers after Run E.}
\end{abstract}

\section{Introduction}
\label{sec:intro}

Fast surrogate simulation of detector effects is a central application of
generative models in high-energy physics~\cite{otus_paper,calogan}.  OTUS
reformulates the problem as an unpaired transport between a physically
interpretable latent space of generator-level particles, $\zspace$, and
detector-level observations, $\xspace$~\cite{otus_paper}.  Its core loss is
\begin{equation}
\label{eq:otus}
\mathcal{L}_{\mathrm{OTUS}}
= \mathbb{E}_{x\sim p(x)}\mathbb{E}_{z\sim p_E(z|x)}
\mathbb{E}_{\tilde x\sim p_D(x|z)}\|x-\tilde x\|^2
+ \lambda\, d_{\mathrm{SW}}\!\left(p_E(z),p(z)\right),
\end{equation}
optionally augmented by directional anchor terms.  The sliced-Wasserstein term
projects high-dimensional particle vectors onto random one-dimensional axes;
our measured failure mode is that a narrow resonance shell is nearly invisible
to most of those axes.

This paper makes two contributions:
\begin{enumerate}
  \item a systematic comparison of the literal OTUS objective, the Run-D
        pair-level extension, and the Virat-style stochastic residual baseline;
  \item the Run-E ``Tier A'' method: physics-aware neural-OT costs,
        learned max-sliced-Wasserstein directions, a cylindrical residual-flow
        parameterization, mass-gated model selection, and modern generative
        diagnostics.
\end{enumerate}

\section{Related Work}
\label{sec:related}

OTUS builds on the sliced-Wasserstein autoencoder~\cite{swae,wae}.  Neural
optimal transport, notably entropic regularization~\cite{cuturi_sinkhorn},
Sinkhorn divergences~\cite{genevay_sinkhorn}, and Monge-gap regularizers,
replaces random projection distances in modern unpaired-transport models.
Flow matching and stochastic interpolants provide continuous-time conditional
generative models~\cite{lipman_flow, albergo_stochastic}; our cylindrical
residual flow is a discrete, physics-constrained step in that direction.
Fast detector simulation has moved from GANs~\cite{calogan} toward
normalizing flows, diffusion models, and equivariant particle-cloud
models~\cite{calochallenge}.  For unfolding, OmniFold is the current
reference~\cite{omnifold}; our encoder remains a single-pass stochastic
unfolding map and is not claimed to replace OmniFold.

\section{Method}
\label{sec:method}

\subsection{Baseline and Run D}
The detector-level data are CMS 2012 DoubleMuParked opposite-sign dimuons in
the $J/\psi$ window.  The Run-D objective extends \ref{eq:otus} with a direct
generator term and a pair-level SWD,
\begin{equation}
\label{eq:rund}
\mathcal{L}_{\mathrm{RunD}}
= \alpha(t)\left[d_{\mathrm{SW}}(z,E(x))
+d_{\mathrm{SW}}(x,D(z))\right]
+ \lambda_{\mathrm{cyc}}\,
\mathbb{E}\|x-D(E(x))\|^2
+ \gamma\,\mathbb{E}\,\mathrm{Huber}\!\left(\frac{m_{\tilde x}-m_x}{m_x}\right).
\end{equation}

\subsection{Run E Tier A components}
\label{sec:rune}

\textbf{Physics-cost Sinkhorn and Monge gap.}
We define a 14-dimensional standardized ground feature vector,
\begin{equation}
u(p)=\bigl[\log\pt^-,\eta^-,\sin\phi^-,\cos\phi^-,
\log\pt^+,\eta^+,\sin\phi^+,\cos\phi^+,
\log m,\log \pt^{\mu\mu},y,\cos\Delta\phi,\sin\Delta\phi,\Delta\eta\bigr],
\end{equation}
and use a weighted squared cost $c(a,b)=\sum_i w_i(u_i(a)-u_i(b))^2$.  The
distribution distance is the debiased Sinkhorn divergence; a Monge-gap proxy
with barycentric projection is available as a regularizer.

\textbf{Max-SW.}  Instead of random slicing directions, a small direction set
is updated adversarially to maximize the sliced distance on each minibatch,
with an orthogonality penalty, and the generative model minimizes the same
distance.

\textbf{Cylindrical residual flow.}  Encoder and decoder are short flows in
cylindrical coordinates.  At each of $K$ steps a physics-conditioned network
predicts a bounded mean residual, a Gaussian core scale, and a Student-$t$ tail
scale; azimuthal coordinates are transported on $S^1$ and energies are
reconstructed on the muon mass shell.

\textbf{Mass-gated selection.}  A checkpoint is eligible only if the generated
mass satisfies hard gates on $W_1$, KS, window fraction, and mean shift.  The
selected-score is set to infinity for ineligible checkpoints.

\textbf{Modern diagnostics.}  We report classifier two-sample test AUC,
fixed-$z$ repeated-decoding stochasticity, and rank-based coverage for cycle
and nearest-neighbor pseudo-pairs.

\section{Experimental Setup}
\label{sec:setup}

Training uses $\approx 580{,}000$ CMS dimuon events and $75{,}904$ MG5 truth
events from the rebuilt ptj$=5$ prior (80/10/10 splits, 20\% CMS cap).  Models
are trained on Apple-silicon MPS.  Evaluation is performed on the untouched
test split.  \todo{Add exact software versions and commit hashes.}

\begin{table}[t]
\centering
\caption{Run D measured test-split mass metrics.}
\label{tab:rund}
\begin{tabular}{lcc}
\toprule
Checkpoint & $W_1$ [GeV] & KS \\
\midrule
three-term baseline & 0.946 & 0.764 \\
Run D best (epoch 270) & 0.026 & 0.221 \\
Run D final (epoch 300) & 0.061 & 0.296 \\
\bottomrule
\end{tabular}
\end{table}

\section{Results}
\label{sec:results}

\subsection{Distribution-level results}
\label{sec:results-dist}
\todo{Insert Run E mass ratio/residual figures from \texttt{plots\_best}.}
Expected reporting table: $W_1$, KS, peak position, width ratio, window
fraction, C2ST AUC, and central coverage for $D(z)$, $E(x)$, and $D(E(x))$.

\subsection{Stochasticity and coverage}
\label{sec:results-stochasticity}
\todo{Insert fixed-$z$ per-event width table and rank-uniformity plots.}

\subsection{Ablations}
\label{sec:ablations}
Planned single-variable ablations: Sinkhorn on/off, max-SW on/off, flow steps
$K=1,2,4$, mass gates on/off.

\section{Discussion}
\label{sec:discussion}

Limitations: no per-event paired CMS truth is available, so conditional
coverage is evaluated only through cycle and nearest-neighbor proxies; the
skim exposes no isolation, impact-parameter, or pileup variables; and the
method is demonstrated for a fixed two-particle final state.  The
generalization to variable particle multiplicity, full event objects, and
Lorentz-equivariant representations remains future work.

\section{Conclusion}
\label{sec:conclusion}
\todo{Final conclusion after Run E.}

\bibliographystyle{unsrtnat}
\bibliography{references}

\end{document}
